
Para extender esta noci\'on de integral a  conjuntos acotados m\'as generales que no sean  paralelepípedos , definimos lo siguiente. 

\begin{definition}\label{def:caja-que-contiene} Sea $W  \subseteq  \mathbb{R}^3$ , $f:W\to\R$ continua y  $P = [a, b] \times [c, d] \times [e, f]$ tal que $W  \subseteq  P$.  Se define la funci\'on $f^*$ tal que
\[
    f^*(x,y,z)=
    \begin{dcases*}
        f(x,y,z) & si $(x,y,z)\in W$ \\[.2cm]
        0        & si $(x,y,z)\notin P\smallsetminus W.$
    \end{dcases*}
\]
La inegral de  $f$ sobre $W$ se define por 
\[
    \iiint_W f\:dV=\iiint_P f^*\:dV.  
\]
\end{definition}

\begin{obs} 
   La \autoref{def:caja-que-contiene} es independiente de la selecci\'on de $P$.
\end{obs}

\begin{propertie}
    Sean $f,\;g$ dos funciones integrables en una regi\'on $W\subset\Rn{3}$, entonces:
    \begin{enumerate}
        \item[i.] $\alpha f+\beta g$ es integrable en $W$, $\forall\;\alpha,\;\beta\in\R$ y adem\'as
        \[
            \iiint_W \left(\alpha f+\beta g\right)dV=\alpha\iiint_W f\:dV+\beta\iiint_W g\:dV.
        \]
        \item[ii.] El producto $fg$ es integrable en $W$.
        \item[iii.] Si $g(x,y,z) \neq  0\;\forall(x,y,z)\in W$ entonces el cociente $f/g$ es integrable en $W$.
        \item[iv.] Si $f(x,y,z) \geq 0  \;\: \forall(x,y,z)\in W$ entonces  $\iiint_W f\:dV\geq0$.
        \item[v.]Si $f(x,y,z) \leq g(x,y,z) \;\: \forall(x,y,z)\in W$ entonces $\iiint_W f\:dV\leq\iiint_W g\:dV.$
        \item[vi.]Si $|f|$ es integrable en $W$, entonces 
        \[
            \left|\iiint_W f\:dV\right|\leq\iiint_W|f|\:dV.  
        \]    
        \item[vii.] Si $W=W_1\cup W_2$  con $W_1\cap W_2 = \varnothing$  (salvo tal vez en conjuntos de medida nula): 
        
        $f$ es integrable en $W$ $\iff$ $f$ es integrable en $W_1$ y $W_2$.
        
         En este caso 
        \[
            \iiint_W f\:dV=\iiint_{W_1} f\:dV+\iiint_{W_2}f\:dV.    
        \]
    \end{enumerate}
\end{propertie}

\begin{definition} Sea $W\subset\Rn{3}$,  se define el volumen de $W$,  y se nota $\text{Vol}(W)$,  a  
    
    \[
        \text{Vol}(W)=\iiint_W 1 \:dV.
    \]
\end{definition}
